\documentclass[aspectratio=169, 14pt]{beamer}
\usepackage{beamer_preamble}

\begin{document}

% ==== Title ===============================================================
\titleslide{Lesson 5-4 \textperiodcentered\ Period 2}{Radical Equations With Rational Exponents + DOK-3 Half-Life Model}

% ==== Do Now ==============================================================
\begin{frame}{Do Now --- Solve the radical equation}{Algebra 2 \textperiodcentered\ Lesson 5-4 \textperiodcentered\ Period 2}
\phasebadges{1}{5}{bridge from P1}
\vspace{0.2cm}
\begin{projcard}{Practice 5-4-savvas-q21 \textperiodcentered\ Cube-root equation}{slidegold}
Solve each radical equation.

\vspace{0.25cm}
\(\sqrt[3]{x} + 8 = 13\)

\vspace{0.25cm}
Use the sentence frame:\quad
``I isolate the radical, then raise both sides to the power \_\_\_.''
\end{projcard}
\end{frame}

% ==== Launch ==============================================================
\begin{frame}{Launch --- Example 4 (Rational-Exponent Equations)}{Algebra 2 \textperiodcentered\ Lesson 5-4 \textperiodcentered\ Period 2}
\phasebadges{2}{12}{teacher models}
\vspace{0.2cm}
\begin{projcard}{Rational-Exponent Rules (keep on your page)}{slidegreen}
\begin{enumerate}
  \setlength{\itemsep}{0pt}
  \item Rewrite \(x^{m/n}\) as \((\sqrt[n]{x})^m\) --- bridge to radicals you know.
  \item To solve \((\text{expr})^{m/n} = k\), raise both sides to reciprocal power \(n/m\).
  \item Simplify, then \textbf{CHECK} --- extraneous solutions arise when denominator is EVEN.
  \item For decay: \(y = a \cdot (0.5)^{(t/h)}\) --- \(h\) is the half-life (denominator of exponent).
\end{enumerate}

\vspace{0.1cm}
\small
\textcolor{slideaccent}{\textbf{EXTRANEOUS SOLUTIONS:}} If \(n\) is EVEN, always substitute back into the original.
\end{projcard}
\end{frame}

% ==== Explore =============================================================
\begin{frame}{Explore --- TPS + DOK-3 Half-Life Model}{Algebra 2 \textperiodcentered\ Lesson 5-4 \textperiodcentered\ Period 2}
\phasebadges{2-3}{33}{student work}
\small\textcolor{slidegray}{5 items in order. Plan \(\sim\)12 min for Practice \#41 (Half-Life Model).}

\vspace{0.15cm}
\begin{projcard}{Try It 4 \textperiodcentered\ Rational-exponent equations}{slideblue}
Solve: \((x^2-3x-6)^{3/2} - 14 = -6\) and \((x-8)^{3/2} = (10-x)^3\).
\end{projcard}
\vspace{-0.15cm}
\begin{projcard}{Practice \#12 \textperiodcentered\ Solve with rational exponent}{slideblue}
\((3x + 2)^{2/3} = 4\)\quad --- does the denominator hint at extraneous solutions?
\end{projcard}
\vspace{-0.15cm}
\begin{projcard}{Practice \#33 \textperiodcentered\ Rational-exponent equation}{slideblue}
\(0.5(x^2+5x+136)^{2/3} = 50\)\quad --- CHECK both solutions.
\end{projcard}
\vspace{-0.15cm}
\begin{projcard}{Practice \#34 \textperiodcentered\ Rational-exponent equation}{slideblue}
\(2(x^2-12x-4)^{1/2} - 3 = 15\)\quad --- CHECK both solutions.
\end{projcard}
\vspace{-0.15cm}
\begin{projcard}{Practice \#41 \textperiodcentered\ PERFORMANCE TASK (\(\bigstar\))}{slideaccent}
Half-Life Model --- identify \(a\), \(h\), \(t\); compute exponent; evaluate; \textbf{INTERPRET in context}.
\end{projcard}
\end{frame}

% ==== Share / Summary =====================================================
\begin{frame}{Share / Summary}{Algebra 2 \textperiodcentered\ Lesson 5-4 \textperiodcentered\ Period 2}
\phasebadges{2}{5}{DOK-3 reveal + P3 bridge}
\vspace{0.2cm}
\begin{projcard}{Sentence frame \textperiodcentered\ Half-Life Model}{slideblue}
``This equation has exponent \_\_\_/\_\_\_, so I raise both sides to the power \_\_\_/\_\_\_.
After simplifying I get \_\_\_.
In the half-life formula, \(a =\) \_\_\_ (start), \(h =\) \_\_\_ (half-life period),
\(t =\) \_\_\_ (elapsed time), giving \(y \approx\) \_\_\_.''
\end{projcard}

\vspace{0.15cm}
\begin{projcard}{Bridge to Period 3}{slideaccent}
Next period: apply rational-exponent skills to \textbf{assessment-critical} items.\\
Bring your check-for-extraneous habit.
\end{projcard}
\end{frame}

% ==== Exit Ticket =========================================================
\begin{frame}{Exit Ticket --- Summary Recap}{Algebra 2 \textperiodcentered\ Lesson 5-4 \textperiodcentered\ Period 2}
\phasebadges{1-2}{3}{not graded}
\vspace{0.2cm}
\begin{projcard}{Complete on your exit ticket}{slidegold}
\begin{enumerate}
  \setlength{\itemsep}{0.2cm}
  \item To solve \((\text{expr})^{m/n} = k\) I raise both sides to power \underline{\hspace{3cm}}.
  \item Half-life tells me the \underline{\hspace{3cm}} of the exponent.
  \item After 16 hours, 50 mL of soft drink decays to \underline{\hspace{2cm}} mL.
  \item One biggest thing I learned today: \underline{\hspace{5cm}}.
\end{enumerate}
\end{projcard}
\end{frame}

\end{document}
